\input regression-test

% Assert the resolved color values rather than comparing typeset output. The
% class is experimental and its box output will change; the palette is what
% this package actually promises.

\makeatletter
\newcommand{\SHOWCOLOR}[1]{%
  \extractcolorspec{#1}\moloch@test@spec
  \typeout{TEST color #1 = \moloch@test@spec}%
}
\makeatother

\ExplSyntaxOn
\NewDocumentCommand \CHECKBODYCOLOR { }
  {
    \color_export:nnN { . } { comma-sep-rgb } \l_tmpa_tl
    \color_export:nnN { mNormaltextFg } { comma-sep-rgb } \l_tmpb_tl
    \tl_if_eq:NNTF \l_tmpa_tl \l_tmpb_tl
      { \typeout{TEST~body~color~matches~the~variant} }
      { \typeout{TEST~body~color~does~not~match~the~variant} }
  }
\ExplSyntaxOff

\begin{document}

\START

\typeout{TEST variant = light (default)}
\CHECKBODYCOLOR
\SHOWCOLOR{mNormaltextFg}
\SHOWCOLOR{mNormaltextBg}
\SHOWCOLOR{mAlertFg}
\SHOWCOLOR{mExampleFg}
\SHOWCOLOR{mFrametitleFg}
\SHOWCOLOR{mFrametitleBg}
\SHOWCOLOR{mProgressbarFg}
\SHOWCOLOR{mTitleseparatorFg}
\SHOWCOLOR{mStandoutFg}
\SHOWCOLOR{mStandoutBg}

% The class's own semantic colors must follow the palette, since that is what
% drives block titles, list labels and contents entries.
\SHOWCOLOR{structure}
\SHOWCOLOR{alert}
\SHOWCOLOR{example}

\molochset{variant = dark}

\typeout{TEST variant = dark}
\CHECKBODYCOLOR
\SHOWCOLOR{mNormaltextFg}
\SHOWCOLOR{mNormaltextBg}
\SHOWCOLOR{mAlertFg}
\SHOWCOLOR{mExampleFg}
\SHOWCOLOR{structure}
\SHOWCOLOR{alert}

% Switching back must restore the selected body color as well as the palette.
\molochset{variant = light}
\typeout{TEST variant = light (restored)}
\CHECKBODYCOLOR

\END

\end{document}
